\section{Step 11: Continuous Abstractness Analysis Using the Question-Level Abstractness Score}
\label{sec:step11_continuous_abstractness}

\subsection{Objective}
The objective of Step~11 is to replace the binary abstract-versus-concrete label with the
\textbf{continuous abstractness score} produced by the question-classification algorithm and test
whether the fNIRS response varies continuously with the level of abstractness.

This step is motivated by the possibility that the binary label may be too coarse.
Two questions may both be labeled \texttt{abstract} while still differing substantially in their
degree of abstractness, and the same is true for questions labeled \texttt{concrete}.
A continuous score preserves this information.

The main question of Step~11 is:
\begin{quote}
In the left frontal ROI, is the duration-aware HbO response associated with the continuous level
of abstractness of the question after adjusting for block structure, question length, and item difficulty?
\end{quote}

\subsection{Why This Step Is Needed}
The previous steps treated abstractness as a binary construct.
That was useful for the first analyses, but it may have discarded meaningful variation.
If the neural effect depends on the \textbf{degree} of abstractness rather than a strict category threshold,
a continuous-score analysis may be more appropriate.

This step therefore treats abstractness as a \textbf{continuum} rather than a dichotomy.

\subsection{Core Analytical Decision}
To keep Step~11 interpretable, the following elements should remain fixed from Step~10:
\begin{enumerate}
    \item the same trial-level frontal outcome;
    \item the same duration-aware variable-window signal summary;
    \item the same included participant-session cohort in the primary analysis;
    \item the same block covariates;
    \item the same item covariates:
    \begin{itemize}
        \item standardized log total word count;
        \item standardized ENEM item correctness.
    \end{itemize}
\end{enumerate}

The main change in Step~11 is:
\begin{quote}
replace the binary condition label with the continuous abstractness score.
\end{quote}

\subsection{Important Modeling Rule}
The binary abstract/concrete label should \textbf{not} be included in the same primary model
as the continuous abstractness score.

The reason is that the binary label is presumably derived from the score by thresholding.
Including both in the same model would create strong collinearity and make the coefficients hard to interpret.

Therefore:
\begin{itemize}
    \item the \textbf{primary model} uses the continuous score only;
    \item the \textbf{binary-label model} is retained only as a benchmark model in a separate analysis.
\end{itemize}

\subsection{Primary Outcome}
The primary outcome remains the same as in Steps~9 and~10:
\begin{quote}
the trial-level left frontal ROI HbO response computed from the duration-aware variable window.
\end{quote}

Keeping the same dependent variable ensures that Step~11 isolates the effect of replacing the
binary label with the continuous abstractness score.

\subsection{Primary ROI Definition}
The primary ROI remains the same left frontal ROI used previously.

\subsubsection{Front ROI HbO channels}
\begin{verbatim}
S2_D6 hbo
S2_D4 hbo
S1_D6 hbo
S1_D3 hbo
S5_D6 hbo
S5_D4 hbo
S5_D3 hbo
\end{verbatim}

\subsubsection{Back ROI HbO benchmark channels}
The back ROI is retained only as a secondary benchmark:
\begin{verbatim}
S4_D2 hbo
S3_D1 hbo
S6_D2 hbo
S6_D1 hbo
S6_D5 hbo
S7_D2 hbo
S7_D5 hbo
S8_D1 hbo
S8_D5 hbo
\end{verbatim}

\subsubsection{Short-separation pairs excluded from inference}
\begin{verbatim}
S1_D8
S2_D9
S3_D10
S4_D11
S5_D12
S6_D13
S7_D14
S8_D15
\end{verbatim}

\subsection{Primary Predictor: Continuous Abstractness Score}
Let \(\texttt{score}_q\) denote the algorithmic abstractness score for question \(q\).

Before fitting the model, the score should be transformed so that:
\begin{quote}
higher values always mean \textbf{more abstract}.
\end{quote}

If the original score is oriented in the opposite direction, it should first be multiplied by \(-1\).

Then define the standardized abstractness score:
\begin{equation}
a_q = z(\texttt{score}_q),
\end{equation}
where \(z(\cdot)\) denotes standardization to mean zero and standard deviation one.

\subsection{Interpretation of the Score Coefficient}
With this coding:
\begin{itemize}
    \item a \textbf{positive} abstractness coefficient means that more abstract questions are associated
    with a larger frontal ROI response;
    \item a \textbf{negative} abstractness coefficient means that more concrete questions are associated
    with a larger frontal ROI response.
\end{itemize}

This sign interpretation should be stated explicitly in the report.

\subsection{Additional Covariates}
Step~11 should retain the same item-level and block-level covariates used in Step~10.

Let:
\begin{equation}
g_{iq} \in \{1,\dots,10\}
\end{equation}
denote realized global block order, and
\begin{equation}
p_{iq} \in \{1,2,3\}
\end{equation}
denote within-block trial position.

Let:
\begin{equation}
w_q = z\!\left(\log(\texttt{total\_word\_count}_q)\right)
\end{equation}
and
\begin{equation}
c_q = z(\texttt{correctness}_q).
\end{equation}

\subsection{Primary Step~11 Model}
Let \(y^{(F)}_{iq}\) denote the duration-aware front ROI HbO response for participant \(i\) on question \(q\).

The primary Step~11 model should be:
\begin{equation}
y^{(F)}_{iq}
=
\beta_0
+
\beta_1 a_q
+
\beta_2 g_{iq}
+
\beta_3 p_{iq}
+
\beta_4 w_q
+
\beta_5 c_q
+
\alpha_i
+
\varepsilon_{iq},
\end{equation}
where:
\begin{itemize}
    \item \(\alpha_i\) is a participant fixed effect;
    \item \(\varepsilon_{iq}\) is the residual term.
\end{itemize}

\subsection{Primary Inferential Target}
The primary inferential target is:
\begin{equation}
\beta_1,
\end{equation}
which represents the association between abstractness level and frontal ROI response after adjustment.

\subsection{Primary Hypothesis}
The primary hypotheses are:

\paragraph{Null hypothesis}
\begin{equation}
H_0: \beta_1 = 0
\end{equation}

\paragraph{Alternative hypothesis}
\begin{equation}
H_1: \beta_1 \neq 0
\end{equation}

The significance threshold should remain:
\begin{verbatim}
alpha = 0.05
\end{verbatim}

\subsection{Primary Inference Strategy}
The preferred inference method should be:
\begin{quote}
two-way cluster-robust standard errors by participant and question
\end{quote}
if the implementation is stable.

If two-way clustered standard errors are unavailable or unstable, the primary fallback should be:
\begin{quote}
participant-clustered standard errors,
\end{quote}
with question-clustered and HC3 standard errors reported as sensitivity analyses.

\subsection{Score Diagnostics Before Modeling}
Before fitting the main model, Step~11 should include a score-diagnostic section.

At minimum, this section should report:
\begin{itemize}
    \item the raw range of the abstractness score;
    \item the mean, standard deviation, median, minimum, and maximum of the score;
    \item a histogram of the score;
    \item the score distribution split by the existing binary abstract/concrete label;
    \item the cutoff used to derive the binary label, if one exists;
    \item the number of questions above and below the cutoff;
    \item the correlation of the score with:
    \begin{itemize}
        \item total word count;
        \item sentence count;
        \item sentence length;
        \item ENEM item correctness;
        \item response time.
    \end{itemize}
\end{itemize}

This section is important because it helps determine whether the abstractness score is mostly capturing
abstractness itself or whether it is strongly entangled with other item properties.

\subsection{Benchmark Model with Binary Label}
To compare the continuous-score approach against the older binary approach, Step~11 should fit one benchmark model:
\begin{equation}
y^{(F)}_{iq}
=
\gamma_0
+
\gamma_1 A_q
+
\gamma_2 g_{iq}
+
\gamma_3 p_{iq}
+
\gamma_4 w_q
+
\gamma_5 c_q
+
\alpha_i
+
\varepsilon_{iq},
\end{equation}
where:
\begin{equation}
A_q =
\begin{cases}
1 & \text{if question } q \text{ is Abstract} \\
0 & \text{if question } q \text{ is Concrete}
\end{cases}
\end{equation}

This model is included only as a benchmark.
Its purpose is to compare:
\begin{itemize}
    \item the continuous-score specification;
    \item the thresholded binary specification.
\end{itemize}

\subsection{Nonlinearity Sensitivity Model}
Because the relationship between abstractness and neural response may not be perfectly linear,
Step~11 should include one quadratic sensitivity model:
\begin{equation}
y^{(F)}_{iq}
=
\beta_0
+
\beta_1 a_q
+
\beta_2 a_q^2
+
\beta_3 g_{iq}
+
\beta_4 p_{iq}
+
\beta_5 w_q
+
\beta_6 c_q
+
\alpha_i
+
\varepsilon_{iq}.
\end{equation}

This model tests whether the relationship is curved rather than strictly linear.

\subsection{Strength-of-Abstractness Sensitivity Model}
If the binary label was derived from a cutoff \(t\), Step~11 may also include a sensitivity model
based on the distance from the cutoff:
\begin{equation}
d_q = z(|\texttt{score}_q - t|).
\end{equation}

A benchmark threshold-strength model may then be:
\begin{equation}
y^{(F)}_{iq}
=
\beta_0
+
\beta_1 A_q
+
\beta_2 d_q
+
\beta_3 g_{iq}
+
\beta_4 p_{iq}
+
\beta_5 w_q
+
\beta_6 c_q
+
\alpha_i
+
\varepsilon_{iq}.
\end{equation}

This is optional, but useful if you want to test whether the neural effect depends on how strongly
abstract or concrete an item is relative to the classification threshold.

\subsection{Secondary Back ROI Benchmark}
Because the back ROI remained null in the previous steps, it should stay secondary here.
Still, one benchmark back-ROI model may be fit:
\begin{equation}
y^{(B)}_{iq}
=
\beta_0
+
\beta_1 a_q
+
\beta_2 g_{iq}
+
\beta_3 p_{iq}
+
\beta_4 w_q
+
\beta_5 c_q
+
\alpha_i
+
\varepsilon_{iq}.
\end{equation}

This model is included only to check whether the continuous score reveals anything in the back ROI
that the binary label did not.

\subsection{Fixed-Window Benchmark Sensitivity}
Because Step~6 used the fixed 7--11~s frontal outcome and because that outcome previously gave one of
the strongest near-threshold frontal results, Step~11 should include one benchmark model using:
\begin{quote}
the Step~6 fixed-window front ROI outcome with the continuous abstractness score.
\end{quote}

This allows a direct comparison between:
\begin{itemize}
    \item the duration-aware continuous-score result;
    \item the fixed-window continuous-score result.
\end{itemize}

\subsection{No-Override Sensitivity}
Step~11 should also include the same no-override sensitivity used in Step~10:
repeat the primary continuous-score model after removing the two inherited manual-inclusion sessions:
\begin{itemize}
    \item PID025 / 2025-11-25 002
    \item PID034 / 2025-11-27 003
\end{itemize}

\subsection{Influence Analysis}
Because the abstractness score is a question-level variable, Step~11 should include a
leave-one-question-out influence analysis.

For each question \(q\):
\begin{itemize}
    \item remove all trials from question \(q\);
    \item refit the primary continuous-score model;
    \item record the new abstractness coefficient;
    \item record the new p-value and confidence interval.
\end{itemize}

This analysis will show whether the apparent abstractness effect depends heavily on one or a few questions.

\subsection{Optional Leave-One-Participant-Out Check}
If computationally feasible, Step~11 may also include a leave-one-participant-out analysis:
\begin{itemize}
    \item remove one participant-session at a time;
    \item refit the primary model;
    \item record the abstractness coefficient and p-value.
\end{itemize}

\subsection{Step-by-Step Workflow}

\subsubsection{Step~11.1: Freeze the Step~10 trial table}
Start from the same Step~10 trial-level duration-aware frontal outcome and the same block metadata.

\subsubsection{Step~11.2: Merge the abstractness score by question}
Add the algorithmic question-level abstractness score to every trial using question ID as the merge key.

\subsubsection{Step~11.3: Align score direction}
Confirm that higher values mean more abstract.
If necessary, reverse the sign of the raw score.

\subsubsection{Step~11.4: Run score diagnostics}
Produce:
\begin{itemize}
    \item histogram of the score;
    \item score distribution by binary label;
    \item score cutoff summary;
    \item correlations with word count, sentence count, sentence length, ENEM correctness, and response time.
\end{itemize}

\subsubsection{Step~11.5: Standardize the continuous predictors}
Compute:
\begin{itemize}
    \item standardized abstractness score;
    \item standardized log total word count;
    \item standardized ENEM correctness.
\end{itemize}

\subsubsection{Step~11.6: Fit the primary continuous-score model}
Fit the participant-fixed-effects model using the continuous abstractness score as the main predictor.

\subsubsection{Step~11.7: Compute robust standard errors}
Use:
\begin{itemize}
    \item two-way cluster-robust SE by participant and question if stable;
    \item otherwise participant-clustered SE as the primary fallback.
\end{itemize}

\subsubsection{Step~11.8: Fit the benchmark binary-label model}
Fit the benchmark model using the thresholded abstract/concrete label instead of the continuous score.

\subsubsection{Step~11.9: Fit the quadratic sensitivity model}
Fit the nonlinear continuous-score model with the squared score term.

\subsubsection{Step~11.10: Fit the optional threshold-strength model}
If a cutoff is available, fit the label-plus-distance benchmark model.

\subsubsection{Step~11.11: Fit the back ROI benchmark}
Fit the secondary back-ROI continuous-score model.

\subsubsection{Step~11.12: Fit the fixed-window benchmark}
Repeat the continuous-score model using the Step~6 fixed-window frontal outcome.

\subsubsection{Step~11.13: Fit the no-override sensitivity}
Repeat the primary continuous-score model after removing PID025 and PID034.

\subsubsection{Step~11.14: Run the leave-one-question-out analysis}
Refit the primary model once per question exclusion.

\subsubsection{Step~11.15: Optionally run the leave-one-participant-out analysis}
Refit the primary model once per participant-session exclusion, if desired.

\subsection{Primary Reporting Quantities}
The primary Step~11 results table must report:
\begin{itemize}
    \item number of included participants;
    \item number of included questions;
    \item number of included trials;
    \item abstractness-score coefficient \(\beta_1\);
    \item standard error;
    \item test statistic;
    \item p-value;
    \item 95\% confidence interval;
    \item block-order coefficient;
    \item within-block-position coefficient;
    \item log-word-count coefficient;
    \item ENEM-correctness coefficient;
    \item robust-SE type used in the primary report.
\end{itemize}

\subsection{Secondary Reporting Quantities}
The secondary Step~11 tables should report:
\begin{itemize}
    \item score-diagnostic summary;
    \item benchmark binary-label results;
    \item quadratic-model results;
    \item optional threshold-strength results;
    \item back-ROI benchmark results;
    \item fixed-window benchmark results;
    \item no-override sensitivity results;
    \item leave-one-question-out coefficient range;
    \item optional leave-one-participant-out coefficient range.
\end{itemize}

\subsection{Primary Conclusion Rule}
The final Step~11 conclusion should be based on the coefficient \(\beta_1\) from the
primary continuous-score model.

The interpretation rule should be:
\begin{itemize}
    \item if \(\beta_1\) is significant and positive, conclude that more abstract questions are associated
    with a larger frontal ROI response;
    \item if \(\beta_1\) is significant and negative, conclude that more concrete questions are associated
    with a larger frontal ROI response;
    \item if \(\beta_1\) is non-significant, conclude that the continuous abstractness score does not
    show a statistically significant linear association with the frontal ROI response in the current dataset.
\end{itemize}

\subsection{Interpretation Templates}

\paragraph{Template if the continuous abstractness effect is significant and positive.}
In the participant-fixed-effects regression using the continuous abstractness score, more abstract
questions were associated with a significantly larger left frontal ROI HbO response
(\(\beta=\ldots\), \(SE=\ldots\), \(p=\ldots\), 95\% CI \([\ldots,\ldots]\)).
This suggests that the neural effect is better described as a continuous abstractness gradient
than as a thresholded abstract-versus-concrete category difference.

\paragraph{Template if the continuous abstractness effect is significant and negative.}
In the participant-fixed-effects regression using the continuous abstractness score, more concrete
questions were associated with a significantly larger left frontal ROI HbO response
(\(\beta=\ldots\), \(SE=\ldots\), \(p=\ldots\), 95\% CI \([\ldots,\ldots]\)).
This suggests that the frontal effect is better characterized as increasing with concreteness rather
than with abstractness.

\paragraph{Template if the continuous abstractness effect is not significant.}
In the participant-fixed-effects regression using the continuous abstractness score, the association
between abstractness level and left frontal ROI HbO response was not statistically significant
(\(\beta=\ldots\), \(SE=\ldots\), \(p=\ldots\), 95\% CI \([\ldots,\ldots]\)).
This suggests that replacing the binary label with a continuous abstractness measure does not by itself
yield robust evidence for a frontal abstractness effect in the current dataset.

\subsection{Expected Outputs}

\subsubsection*{Tables}
\begin{verbatim}
01_step11_score_diagnostics.csv
02_step11_score_correlations.csv
03_step11_primary_continuous_model.csv
04_step11_binary_benchmark_model.csv
05_step11_quadratic_model.csv
06_step11_threshold_strength_model.csv
07_step11_back_roi_benchmark.csv
08_step11_fixed_window_benchmark.csv
09_step11_no_override_results.csv
10_step11_leave_one_question_out.csv
11_step11_leave_one_participant_out.csv
12_step11_final_conclusion.csv
\end{verbatim}

\subsubsection*{Figures}
\begin{verbatim}
figures/step11/score_histogram.png
figures/step11/score_by_binary_label.png
figures/step11/score_vs_wordcount.png
figures/step11/score_vs_correctness.png
figures/step11/continuous_partial_effect_plot.png
figures/step11/continuous_vs_binary_benchmark.png
figures/step11/quadratic_effect_plot.png
figures/step11/no_override_comparison.png
figures/step11/leave_one_question_out_plot.png
\end{verbatim}

\subsubsection*{Reports}
\begin{verbatim}
reports/step11/step11_continuous_abstractness_report.tex
reports/step11/tables/primary_continuous_table.tex
reports/step11/tables/benchmark_table.tex
reports/step11/text/final_conclusion.tex
\end{verbatim}

\subsection{Minimal Figures for the Main Report}
The main Step~11 report should include at least:
\begin{enumerate}
    \item histogram of the abstractness score;
    \item score distribution split by the binary abstract/concrete label;
    \item scatter plot of score versus word count;
    \item scatter plot of score versus ENEM correctness;
    \item partial-effect plot for the primary continuous-score model;
    \item benchmark comparison between the continuous-score and binary-label models;
    \item leave-one-question-out abstractness coefficient plot.
\end{enumerate}

\subsection{Quality-Control Checks}
Before Step~11 is considered complete, verify that:
\begin{enumerate}
    \item the Step~10 frontal outcome is unchanged;
    \item the abstractness score is merged correctly by question ID;
    \item the direction of the score is aligned so that higher means more abstract;
    \item the continuous score is not included in the same primary model as the binary label;
    \item the primary robust-SE specification is documented clearly;
    \item the no-override sensitivity removes only PID025 and PID034;
    \item leave-one-question-out runs cover all included questions;
    \item the score diagnostics are reported before interpretation.
\end{enumerate}

\subsection{Minimal Completion Criteria}
Step~11 is complete when:
\begin{itemize}
    \item the abstractness score has been merged and validated;
    \item the score-diagnostic section has been completed;
    \item the primary continuous-score model has been fit;
    \item the binary benchmark model has been fit;
    \item the quadratic sensitivity model has been fit;
    \item the back-ROI benchmark has been completed;
    \item the fixed-window and no-override sensitivities have been completed;
    \item the leave-one-question-out analysis has been completed;
    \item the final conclusion has been written.
\end{itemize}

\subsection{Short Practical Summary}
In simple terms, Step~11 should:
\begin{quote}
take the same frontal duration-aware outcome from Step~10, replace the binary abstract/concrete
label with the continuous abstractness score from the classification algorithm, test whether the
frontal signal changes continuously with abstractness level, compare that result against the old
binary-label benchmark, and determine whether the neural effect is better described as a gradient
than as a thresholded category difference.
\end{quote}